\section{嵌套虚拟化：从Turtles到HyperTurtle}\label{sec:nested}

前两章分别考察了纯软件模拟与以 Arm 为代表的硬件辅助虚拟化，二者处理的都是"一层 hypervisor 直接管理一层 guest"的问题。本章把对象再推进一层：当 hypervisor 自己就跑在一台虚拟机里时，机器状态、地址映射与事件策略要如何组合才既正确又可接受地快。本章的架构口径限定为 x86-64，即 amd64 这一 ISA 称谓，讨论的机制都落在 Intel 的 VMX 指令集、VMCS 控制结构与 EPT 扩展页表上，不使用 AMD SVM 的 VMCB/NPT 术语，也不把 Arm 的 EL2 异常级名称移植过来。两篇核心文献分别位于这条问题链的两端：OSDI 2010 年的 Turtles 回答了"怎样在缺乏架构支持时正确实现嵌套"~\cite{source7}，USENIX ATC 2025 年的 HyperTurtle 则回答了"在已经能嵌套的前提下怎样减少不必要的跨层往返"~\cite{source8}。前一问是后一问的前提，二者不能按同一标准横向比较性能，只能按"谁控制资源、关键路径怎么走、何时介入、代价多大"这条主线先后展开。

\subsection{问题建立：当hypervisor自己也在虚拟机里}\label{sec:nested-problem}

嵌套虚拟化的层级用下标约定：L0 是跑在裸机上的 hypervisor，L1 是运行在 L0 虚拟机里的 guest hypervisor，L2 则是 L1 自己管理的嵌套 guest。L1 得到的机器接口必须包含它管理 L2 所需的虚拟化能力，即使真实 VMX 资源仍由 L0 控制；L2 则使用 L1 向它提供的虚拟机器接口~\cite{source7,source8}。

这一层次关系来自具体的部署需求。最直接的一类场景是租户自治：云厂商在 L0 上管理物理机，租户却希望在租给自己的虚拟机里再起 hypervisor，以便自行决定里面跑什么、如何迁移、何时快照，而无需惊动基础设施方。第二类是虚拟化软件自身的测试与开发——要复现"hypervisor 里再跑 hypervisor"的复杂环境，单层虚拟化天然不够。第三类，也是 HyperTurtle 论文最关心的，是基于虚拟机的安全容器：Kata 这类运行时用 microVM 而不是共享内核的进程来承载容器，换取更强的隔离边界；当这一层容器运行时本身又跑在租户的 VM（L1）里时，microVM 就成了 L2，嵌套就不再是可选项而是部署形态本身。必须澄清的是，并非所有容器都需要嵌套——普通容器共享宿主内核，不构成 L2；只有"在虚拟机内部再起 microVM"这一形态才触发本节讨论的问题，因此不能把容器隔离一律等同于嵌套虚拟化~\cite{source8}。

问题的难点不在"再多一层特权环"这样的直觉，而在资源归属的不对称。L1 确实要使用虚拟化扩展：它要执行 VMX 指令、创建自己的 VMCS、配置自己的 EPT。但真正的硬件资源——处理器当前运行所需的 VMX 状态、活动 EPT 上下文与真实 IOMMU——都由 L0 独占控制。于是 L0 面临的任务不是给 L1 一个更强的 ring，而是向 L1 提供一个"虚拟的虚拟化接口"：把 L1 要求的虚拟化操作接住、与 L0 自身对硬件和其余租户的约束合并，再组织成处理器当前可执行的一组硬件状态~\cite{source7}。

代价的直观体现是事件路径被拉长。当 L2 里发生一个需要 L1 处理的 VM-exit 时，CPU 先退出到 L0，L0 判断该事件需要向 L1 报告，并向 L1 反射退出，L1 处理完毕请求恢复 L2，L0 再次介入执行 VMEntry，最后才回到 L2，形成 $L2 \to L0 \to L1 \to L0 \to L2$ 的往返。但这条链不是所有事件都要走完：外部中断、NMI，以及只有 L0 在 VMCS02 里设置了 trap 而 L1 没在 VMCS12 里设置的事件，属于 L0 自己的范畴，L0 就地处理后直接恢复 L2，从 L1 的视角看这些事件根本不存在——正如单层虚拟化里 CPU 透明处理一次硬件事件后 hypervisor 继续运行一样~\cite{source7}。事件归属的判定权始终在 L0 手里，这才是"L0 保留最终控制"的准确含义，而不是"所有事件逐层转发"。

\subsection{Turtles：把多层语义组合到单层硬件支持上}\label{sec:turtles}

Turtles 是 KVM 之上的嵌套虚拟化实现，目标是运行未经修改的 guest hypervisor（KVM、VMware Server 等）及其 guest 系统。它在当时缺少专门嵌套硬件支持的条件下实现了可用的嵌套系统，并通过多类负载测量开销；论文的低开销结果不是对所有负载的统一上界。其实现涉及：CPU 上的 VMCS 合成、内存上的多维分页、I/O 上的多级设备分配，以及贯穿其中的退出归属判定~\cite{source7}。下面按四个维度分别说明"谁读写什么、怎样合成、代价在何处"。

\subsubsection{CPU：虚拟化VMX与合成VMCS}\label{sec:nested-cpu}

在论文使用的单层 VMX 基线上，L1 在 VMX non-root 环境执行虚拟化管理指令会触发相应的拦截，这就给了 L0 一个天然的介入点：L0 捕获并模拟 L1 的 VMX 指令，也就是"trap-and-emulate"。大部分 VMX 指令的处理是"L1 退出到 L0、L0 按 VMCS 更新规则改写状态、再回到 L1"，让 L1 觉得指令真的执行了；但用来启动虚拟机的 \code{vmresume} 与 \code{vmlaunch} 是例外，因为此时 L0 必须模拟一次 L1→L2 的 VMEntry，流程变成"L1 退出到 L0，L0 直接 VMEntry 进入 L2"~\cite{source7}。

状态合成的关键是三份 VMCS 的分工，如图~\ref{fig:nested} 所示。VMCS01 是 L0 运行 L1 用的结构，从硬件角度看它和单层虚拟化完全一样：一个 hypervisor（L0）在 root mode 跑，一个 guest（L1）在 guest mode 跑。L1 为 L2 建立 VMCS12 来描述所需运行环境。该结构的地址与 host 状态具有 L1 视角下的含义，L0 不能不经检查、转换和合成就把它作为真实控制结构装载。于是 L0 以 VMCS12 为蓝本，再叠加 VMCS01 里对 L1 的约束，合成 VMCS02 作为真正运行 L2 的结构：guest state 大多可直接拷贝，host state 必须填成 L0 自己的、以便 CPU 从 L2 正确切回 L0，而 VMCS12 的 host state 要拷进 VMCS01 的 guest state，使 L0 模拟 L2→L1 切换时能加载正确的 L1 上下文~\cite{source7}。

控制数据的合成是正确性的重心。设想 L1 在 VMCS12 里要求 trap 某事件 $E_A$，但 L0 在 VMCS01 里并没有为 L1 trap 这个事件，那么为了把 $E_A$ 转发给 L1，L0 必须在 VMCS02 里补上对应的 trap；同理，L1 用来向 L2 注入事件的字段、以及 CPU 用来报告退出原因的字段都要合并。这套合并的具体实现落在 KVM 源码里，论文没有逐字段展开~\cite{source7}。这里有一个必须澄清的术语边界：Turtles 所说的 VMCS shadowing 是 L0 在软件层面"影子维护"VMCS12 并合成 VMCS02 的技术，与后来 Intel 在 Haswell 引入的同名硬件特性（让 L0 不必在每次 \code{vmread}/\code{vmwrite} 时都 trap）不是一回事；Turtles 的机制不依赖任何后续硬件能力~\cite{source7,source8}。同样，Turtles 描述的 2010 年合成逻辑，不应被当成当前 KVM 全部路径的写照——如今内核里的实现（如 \path{arch/x86/kvm/vmx/nested.c} 的 \code{prepare_vmcs02}）已经历多次重构并叠加了硬件加速，只可作为"合成 VMCS02 这一思路延续至今"的旁证，而非当年原型的源码~\cite{source5}。

\subsubsection{内存：多维地址映射与合成EPT}\label{sec:nested-mem}

地址翻译在嵌套下多出一层：逻辑链是 L2 VA→L2 GPA→L1 GPA→host PA，即三层表。但当时的硬件 MMU 最多只支持两维（guest 页表 + EPT/NPT），无法同时走三层。Turtles 的做法是把三层折叠到硬件提供的两层上：L2 的常规页表保留下来负责 L2 VA→L2 GPA 这一层，L0 则合成一张 EPT02 直接完成 L2 GPA→host PA，而不是要求硬件逐个遍历 L2 页表、EPT12、EPT01~\cite{source7}。图~\ref{fig:nested} 的上下两行分别给出逻辑映射与实际硬件执行路径，其中 EPT 下标"管理层、被管理层"表示该表由谁为谁维护，不表示地址转换方向。

这张合成 EPT02 是按需建立的。L0 以一张空 EPT02 启动 L2，当 L2 访问的 GPA 在 EPT02 里缺项时触发 EPT violation 退出到 L0；L0 检查 L1 维护的 EPT12 是否已有该地址的合法映射——有，则把 EPT12 里的 L2→L1 映射与 EPT01 里的 L1→host 映射拼接，写入 EPT02 的对应项；没有，说明这个缺页属于 L1 应负责的范围（例如 L1 尚未为该 GPA 分配内存），L0 就把退出交给 L1 处理，待 L1 更新 EPT12 后再回来补 EPT02~\cite{source7}。为保持 EPT02 与 EPT12 一致，L0 把 EPT12 所在内存设为只读，L1 一旦改写就 trap，L0 据此同步 EPT02；同时 L0 还要 trap L1 的 \code{INVEPT} 以相应失效 EPT 缓存。这套机制相比两个朴素方案的优势在图~\ref{fig:nested} 上体现为：shadow-on-shadow 要软件维护两层影子页表并压缩成 SPT02，代价最高；shadow-on-EPT 让 L1 用影子页表跑 L2，虽然用上了 EPT 硬件，但 L2 的每次页故障都要 L1 处理、每次页表写都要 L0 转发，等于把单层影子页表慢的原因在嵌套里又放大一遍。多维分页减少了影子页表维护引起的额外退出。EPT02 缺项或权限问题等仍可导致 EPT violation；L2 自身的普通页故障通常由 L2 处理，除非上层配置了相应拦截~\cite{source7}。此外 L0 还允许 L1 使用 VPID 给 TLB 条目打标签以避免每次切换都冲刷 TLB，并负责把 L1 自选的 VPID 重映射到 L0 的有效编号，防止两层各自随意取值时冲突；用大页给 guest 内存做后备还能进一步缩小 EPT 表~\cite{source7}。

\subsubsection{中断：事件归属、转发与注入}\label{sec:nested-irq}

嵌套下的中断必须分清三件不同的事：宿主自己产生的外部中断或 NMI，L0 要注入给 L1 的虚拟中断，以及 L1 希望注入给 L2 的事件。它们的处理路径不同。当 L2 正在运行且发生一次退出时，L0 首先按 VMCS12 与 VMCS02 的 trap 配置判断归属：属于 L0 的事件（外部中断、NMI，或只在 VMCS02 里 trap 的事件）就地处理后恢复 L2；属于 L1 的事件，L0 把 CPU 已更新的 VMCS02 字段拷回 VMCS12，再恢复 L1，让 L1 觉得是一次 L2→L1 的直接退出~\cite{source7}。注入方向上，L1 在 VMCS12 里写的注入字段要被 L0 合并进 VMCS02，才能让 CPU 真正把中断送进 L2。

设备分配场景揭示了一个关键的非对称性：DMA 直达不等于中断直达。多级设备分配让 L2 的数据可以直接落到设备，但中断在当时的 x86 架构上无法直接派发给 guest——中断到来总是先退出到 L0，L0 再转发 L1，L1 注入 L2，其中断的 EOI 写回也会再次触发退出。论文据此推断，设备中断及其 EOI 造成的退出是设备分配场景下残余开销的最大来源，并用一个去掉中断、纯轮询的网络驱动做了对照实验来支撑这一假设~\cite{source7}。需要强调的是，这是 2010 年原型在当时架构下的路径描述；后续硬件与 KVM 实现继续演进，因而不能把当年的中断处理路径断言为所有当前实现的必然行为~\cite{source7,source8}。

\subsubsection{设备：多级设备分配与虚拟IOMMU}\label{sec:nested-dev}

两层虚拟化各可选用模拟、半虚拟化或设备分配三种 I/O 方式，组合出九种可能，Turtles 重点评估了其中五种，并把"L0→L1 与 L1→L2 都做设备分配"称为多级设备分配~\cite{source7}。其目标不是把真实网卡"直通"给 L2 就完事，而是要在数据路径外补全三条辅助路径。

DMA 的地址问题是第一道坎：guest 驱动用的是 guest 物理地址，设备访存用的是 host 物理地址，而当时的 IOMMU 只支持一级翻译，两层映射必须再次压缩到一层。对可修改的 L1 可用半虚拟化 IOMMU——把 L1 里设置 L2→L1 映射的代码替换成对 L0 的 hypercall，L0 把映射里的 L1 地址换成对应 L0 地址后写入真实 IOMMU；对不改 L1 的方案，L0 为 L1 模拟一个 IOMMU，L1 以为自己在真 IOMMU 上建 L2→L1 映射，L0 截获这些映射、重映射 L1 地址、把 L2→host 的映射写到真实 IOMMU 上~\cite{source7}。中断路径如前所述仍需退出转发。MMIO 与 PIO 则相对干净：嵌套 guest 的 MMIO/PIO 与单层 guest 走同样的机制，不在 I/O 关键路径上引入额外退出~\cite{source7}。也就是说，多级设备分配的价值在于把 DMA 与 MMIO/PIO 留在直通路径上，代价集中在中断转发与 L1 注入这两类退出上，隔离和通知并没有被"直通"二字抹掉。

\begin{figure}[tbp]
\centering
\begin{tikzpicture}[x=1cm,y=1cm]
\node[box,text width=2.6cm] (v01) at (0,0) {VMCS01\\L0 运行 L1};
\node[box,text width=2.6cm] (v12) at (4.6,0) {VMCS12\\L1 描述 L2};
\node[greenbox,text width=2.6cm] (v02) at (9.2,0) {VMCS02\\L0 实际运行 L2};
\draw[wire] (v12)--node[above,font=\footnotesize]{状态与控制合成}(v02);
\draw[dwire] (v01.south)--++(0,-0.6)-|node[pos=0.4,below,font=\footnotesize]{同时满足 L0 的控制要求}(v02.south);
\node[note,anchor=west] at (-1.3,-1.45) {逻辑地址映射};
\node[box,text width=1.9cm] (lva) at (0,-2.2) {L2 VA};
\node[box,text width=1.9cm] (lgpa) at (3.4,-2.2) {L2 GPA};
\node[box,text width=1.9cm] (l1pa) at (6.8,-2.2) {L1 GPA};
\node[box,text width=1.9cm] (hpa) at (10.2,-2.2) {host PA};
\draw[wire] (lva)--node[above,font=\footnotesize]{L2 页表}(lgpa);
\draw[wire] (lgpa)--node[above,font=\footnotesize]{EPT12}(l1pa);
\draw[wire] (l1pa)--node[above,font=\footnotesize]{EPT01}(hpa);
\node[note,anchor=west] at (-1.3,-3.1) {实际硬件执行路径};
\node[greenbox,text width=1.9cm] (ava) at (0,-3.9) {L2 VA};
\node[greenbox,text width=1.9cm] (agpa) at (3.4,-3.9) {L2 GPA};
\node[greenbox,text width=1.9cm] (ahpa) at (10.2,-3.9) {host PA};
\draw[wire] (ava)--node[above,font=\footnotesize]{L2 页表}(agpa);
\draw[wire] (agpa)--node[above,font=\footnotesize]{合成 EPT02}(ahpa);
\end{tikzpicture}
\caption[Turtles 的控制与地址映射合成]{Turtles 的控制与地址映射合成。VMCS/EPT 下标表示管理层和被管理层；EPT02 由 L0 按需建立并维护一致性。\cite{source7}}
\label{fig:nested}
\end{figure}


图~\ref{fig:nested} 把 CPU 与内存两个维度的合成放在同一张图里：上排是 VMCS01/12/02 的控制合成关系，下排是地址映射的逻辑链与硬件执行链的对照。该图同时说明，Turtles 的"组合"本质上是把 L1 的语义要求（VMCS12、EPT12）作为输入，与 L0 自身的约束（VMCS01、EPT01）合并成一份硬件可执行的产物（VMCS02、EPT02），而合并的依据是退出归属与地址归属的判定规则。

\subsection{从"能嵌套"到"少往返"：性能问题的重新定位}\label{sec:nested-cost}

Turtles 的成就是正确性与可行性：在不修改 guest hypervisor 的前提下给出了嵌套语义。但它的开销分析也暴露了问题的下一层形态。论文用 L2 里循环执行 \code{cpuid} 的合成微基准测单次退出的周期成本：裸机上 \code{cpuid} 约 100 周期，单层虚拟化约 2600 周期（其中约 1000 周期是模式切换），嵌套则飙到约 58000 周期~\cite{source7}。原因是 L1 每处理一次 L2 退出，都要执行多条 \code{vmread}/\code{vmwrite} 这类特权指令，而每条又各自触发一次对 L0 的退出——退出被逐级放大。论文在这里区分两类优化。L0 的合并路径可以只拷贝变化字段，并探索直接批量访问真实 VMCS 的方式；后一尝试不严格符合 VMX 规范，不能保证跨处理器适用。另一个称为 DRW 的评测配置则明确修改 L1，使其直接读写 VMCS12，用来估计省去 L1 的 VMREAD/VMWRITE 拦截的潜在收益。这两类直接访问不能混称为只修改 L0 的同一个优化~\cite{source7}。

HyperTurtle 的动机分析把这个现象测得更系统，给出两条明确的量化结论。其一，从 VM-exit 到下一次 VM-entry 的处理时长因退出原因而异，其中包含切换和软件处理，不能将整个数值视为纯硬件切换延迟：L0 处理一次 nop hypercall 从 L1 退出约 0.85$\mu$s，从 L2 退出则约 2.82$\mu$s；NMI 的差距更大，从 L1 是 5.7$\mu$s，从 L2 是 16.7$\mu$s~\cite{source8}。其二，L2 触发的退出普遍比 L1 更贵，因为 L0 需要把退出注入给 L1、并维护合成控制结构及其他管理状态；论文援引 Vilanova 等人的结果指出，world-switch 相关开销约占一次 10$\mu$s 的 \code{cpuid} 模拟的 33\%~\cite{source8}。这两个结论共同指向：跨层往返及其处理是这些负载的重要成本来源，不能只按地址翻译层数解释嵌套开销。

于是出现一个张力。想绕过 L1 的优化确实存在——DVH（Direct Virtual Hardware）把 L0 的虚拟设备直接分配给 L2，从 L2 的数据路径上拿掉 L1，延迟与吞吐都明显改善。但代价是 L1 对 L2 失去了控制：L1 不再处于 L2 的网络控制路径上，无法再施加路由、监控、限速等策略，而这些正是容器管理所必需的~\cite{source8}。"路径更短"与"管理权仍在租户"由此对立起来。HyperTurtle 试图回答的过渡问题因此不是"怎样把 L1 整个搬进 L0"，而是：能否让一部分由 L1 定义的处理逻辑在 L0 就地执行，既省掉往返、又保留 L1 的策略控制，同时不破坏 L0 的资源边界与隔离。

\subsection{HyperTurtle：把L1的部分处理逻辑下放到L0}\label{sec:hyperturtle}

HyperTurtle 的答案是把 L1 hypervisor的一部分功能以 eBPF 程序的形式下放到 L0，称之为 hyperupcall：L1 把自己对某类 L2 事件的处理逻辑写成一个 eBPF 程序，交给 L0 在 L2 事件的处理路径上就地调用，从而在"L1 仍然定义逻辑、L0 负责执行与兜底"的前提下省去进入 L1 的往返。其设计目标有四条：消除多余 world-switch（性能）、保留 L1 对 L2 的策略控制（控制）、最小化对 L0 的改动（低侵入）、提供可推广到不同子系统的通用机制（通用）~\cite{source8}。实现面向 KVM，且限定为单层嵌套。

\subsubsection{通用机制：注册、共享状态、执行与回退}\label{sec:nested-ht-infra}

注册路径与执行热路径是两回事，必须分开看。安装由 L1 通过 hypercall 发起，经 L0 的 KVM 与 QEMU 等用户态组件落地：L1 提供编译好的 eBPF 程序与要挂接的 hook，L0 借 Linux 内核已有的 eBPF 注册机制完成验证、加载与链接；与更早的 hyperupcall 原型不同，HyperTurtle 把"加载"与"链接"拆成两个 hypercall，使同一个 hyperupcall 能挂到多个 hook 而无需重复加载~\cite{source8}。安装完成后，事件热路径可调用已链接的 eBPF：例如 L2 退出到 L0 后，就地完成可下放的处理，再恢复 L2；网络 hook 则位于相应的数据处理路径。图~\ref{fig:hyperupcall} 对比的是普通处理与下放后的事件路径，安装动作不在图示热路径之内。

下放逻辑要在 L0 里正确运行，需要两类共享状态。一是数据通道：L0 与 L1 之间通过共享的 eBPF map 交换数据，而不是早先 hyperupcall 的基于数据包的方式，以支持并发访问；具体做法是新增一个 hypercall 创建一个共享内存 PCI 设备，把数组型 eBPF map 暴露给 L1。二是受限的内存访问：新增 helper \code{bpf_probe_read_hyperupcall} 让 eBPF 能读 L1 的物理内存，该 helper 会校验地址有效且确属发起该程序的 L1，从而把"每个 hyperupcall 只能读自己所属 L1 的内存"这一隔离写成强制约束~\cite{source8}。

程序验证只是信任边界的一部分，访问授权与状态同步同样需要单独设计。HyperTurtle 支持两种安全模型：一是 L1 提交的程序由 L0 的 eBPF verifier（或更严格的 Prevail 之类）验证；二是只接受云厂商预先审核、签名验证过的程序。无论哪种，隔离都由 verifier、受限 helper 的地址绑定、以及共享状态的同步协议三者联合保证，任何单独一环都不足以宣称跨租户安全~\cite{source8}。

最后，下放永远带着回退。HyperTurtle 明确保留了 L1 的原始处理路径作为兜底：内存池耗尽、fault log 写满、加锁失败、遇到不支持的情况时，一律回退到传统的、需要进入 L1 的处理方式~\cite{source8}。这意味着 hyperupcall 优化的是"适用热路径"而不是取消 L1；L1 仍是功能完整性与复杂功能的最终承担者。

\subsubsection{主案例：EPT fault处理}\label{sec:nested-ht-ept}

EPT fault 展示了同一逻辑操作引发多次跨层往返的情形。论文的 EPT hyperupcall 面向 shmem 文件后备内存，不能直接当作所有内存后备方式的通用处理器。图~\ref{fig:hyperupcall} 已给出通用对比，这里补上具体的步数差异：L2 因 EPT02 缺项触发退出到 L0（第 1 步），L0 检查 EPT12 是否有合法映射（第 2 步）；若没有，L0 把 EPT fault 注入 L1 并进入 L1（第 3 步），L1 在 EPT12 里建立映射（第 4 步），L1 试图恢复 L2、又 trap 回 L0（第 5 步），L0 进入 L2（第 6 步），L2 重执行触发 fault 的指令——由于 EPT02 仍缺项，再次退出（第 7 步），L0 再次查 EPT12（第 8 步）、找到条目、用其补全 EPT02，最后再进入 L2（第 9 步）。一次缺页共产生 6 次 world-switch，而非嵌套下的 2 次~\cite{source8}。EPT fault hyperupcall 做的事，就是把这个"缺项→进 L1 建映射→重执行→再缺项→再查表"的循环里属于 L1 的那段，替换成 L0 内的一段 eBPF：在发现 EPT12 缺项时，hyperupcall 先检查已有映射和并发状态；适用时复用映射或从预分配内存池取帧，再更新 EPT12 与 fault log，L0 随即完成 fault 处理并恢复 L2，省掉了进入 L1 及其后的重执行往返~\cite{source8}。

正确性的难点全部在共享状态上。第一是内存来源：eBPF 上下文里分配内存过于复杂，因此 L1 维护一个预分配的帧池放在共享 map 里，hyperupcall 从池里取帧，池空则回退；实验表明 4096 帧足以避免同步补充，论文全文用这一池大小~\cite{source8}。第二是与 L1 的一致性：hyperupcall 改动了 EPT12，L1 必须知道；它通过共享的环形 fault log 记录自己映射过的 L2-PA，L1 在下一次进入 L1 的 vm-exit 时把这些条目对应的 L1-VA 映射进 VMM 地址空间；log 写满则回退。反过来，VMM（或 virtiofsd 等外部进程）可能先行建了映射而 hyperupcall 尚未察觉，因此 L1 把 VMM 的 cr3 拷进共享 map，让 hyperupcall 能用 \code{bpf_probe_read_hyperupcall} 遍历 VMM 页表、判断该地址是否已有映射；同时把 shmem 的页故障处理函数包一层，先加锁、再查 EPT12 是否已被 hyperupcall 映射、再走原处理、最后把对应的 L2-PA 写进环形 page fault log，以弥合"hyperupcall 已建映射但 L1 尚未看到"的窗口~\cite{source8}。第三是并发：整个 hyperupcall 包在一把锁里，锁也经共享 map 暴露给 L1，使 L1 处理 EPT fault 时能取同一把锁；受 eBPF 锁实现限制，取锁失败即回退~\cite{source8}。核心评议因此不在"eBPF 跑得快"，而在"减少跨层切换之后，如何仍正确维护两层内存管理状态"——池、双日志、cr3 共享与锁，本质上是把原本主要由 L1 内存管理路径维护的状态关系，改造成 L0 与 L1 之间可并发、可回退的协议。

\subsubsection{扩展案例：网络策略与性能采样}\label{sec:nested-ht-net-prof}

网络案例回应的是 DVH 的取舍。DVH 把 L0 提供的虚拟网卡直接分配给 L2，数据路径短、性能好，但 L1 失去策略控制；HyperTurtle 的做法不是退回两层 virtio，而是让 L1 把过滤、限速、监控逻辑写成 eBPF 挂到 L0 的网络接口 hook 上，从而在保留 DVH 直通数据路径的同时恢复 L1 的控制。论文实现了四个示例程序：什么都不做的 pass、按源 IP 过滤的无状态防火墙、可运行时调参的限速器、以及监控 L2 全部 TCP 连接的 TCP-Top~\cite{source8}。必须澄清的是，这里的"分配"是把 L0 的虚拟网卡（vNIC）交给 L2，而不是物理网卡直通；隔离靠"每个 hyperupcall 只能跑在它所属 L1 拥有的虚拟接口上"来保证，且每个 vNIC 可挂不同程序，使不同 L2 各受其 L1 策略约束。为接回 Kubernetes，论文实现了一个 CNI 插件做动态接口管理：pod 创建时经 hypercall 请求 L0 建 vNIC、由 QEMU 热插到 L1、CNI 再配置网络命名空间与地址，从而在 k8s 管理面下既用 DVH 又保留 L1 的网络函数~\cite{source8}。这个案例的意义在于展示通用性：hyperupcall 不局限于内存，还能把"数据处理"与"策略控制"重新拼在同一层。

profiling 案例则展示另一种下放形态：采样本身不需要 L1 提供信息，只需把样本送回 L1。普通情况下，从 L1 对 L2 采样一次，L0 注入 NMI、L1 处理采样并访问 MSR、\code{vmread}/\code{vmwrite} 改 VMCS、再执行 \code{vm-enter}，全程约 26 次 world-switch；其中 L0 注入 NMI 到 L1 要 7.8$\mu$s，模拟 L1 的 \code{vm-enter} 还要 8$\mu$s~\cite{source8}。HyperTurtle 借 Linux 已有的 perf 接口把 eBPF 注册成采样事件，新增 helper \code{vcpu_probe} 让 hyperupcall 直接读 L2 的寄存器，样本经共享 map 的环形缓冲送回 L1；注册过程把采样限制到相应执行对象的 guest 状态，并结合受控寄存器读取接口维持访问边界~\cite{source8}。两个案例合起来说明方法通用，但结论有边界：它们证明的是"特定处理路径可以减少往返"，而不是"所有中断都能直接注入"或"所有设备都能完整下放"。

\begin{figure}[tbp]
\centering
\begin{tikzpicture}[x=1cm,y=1cm]
\node[note,anchor=west] at (-1.0,0.9) {普通嵌套：以一个需要 L1 处理的退出事件为例};
\node[box,text width=1.5cm] (a) at (0,0) {L2};
\node[box,text width=1.5cm] (b) at (3,0) {L0};
\node[amberbox,text width=1.5cm] (c) at (6,0) {L1\\处理事件};
\node[box,text width=1.5cm] (d) at (9,0) {L0};
\node[box,text width=1.5cm] (e) at (12,0) {L2};
\draw[wire] (a)--node[above,font=\footnotesize]{退出}(b);
\draw[wire] (b)--node[above,font=\footnotesize]{反射}(c);
\draw[wire] (c)--node[above,font=\footnotesize]{请求恢复}(d);
\draw[wire] (d)--node[above,font=\footnotesize]{进入}(e);
\node[note,anchor=west] at (-1.0,-1.15) {HyperTurtle：已注册的处理逻辑在 L0 执行};
\node[greenbox,text width=1.5cm] (f) at (0,-2.1) {L2};
\node[greenbox,text width=5.0cm] (g) at (6,-2.1) {L0：调用 L1 的 eBPF hyperupcall\\共享 maps、受限 helper、状态同步};
\node[greenbox,text width=1.5cm] (h) at (12,-2.1) {L2};
\draw[wire] (f)--node[above,font=\footnotesize]{退出}(g);
\draw[wire] (g)--node[above,font=\footnotesize]{成功后恢复}(h);
\node[amberbox,text width=3.1cm] (fallback) at (6,-4.0) {L1 原有处理路径\\保留正确性与复杂功能};
\draw[dwire] (g)--node[right,font=\footnotesize]{资源不足 / 竞争 / 不支持时回退}(fallback);
\end{tikzpicture}
\caption[HyperTurtle 减少跨层往返]{HyperTurtle 减少跨层往返的方式。图显示通用热路径，具体 EPT 缺页还可能包含重执行后再次退出等步骤；成功下放只减少特定处理所需的切换，不取消 L1。\cite{source8}}
\label{fig:hyperupcall}
\end{figure}


图~\ref{fig:hyperupcall} 概括了 HyperTurtle 的通用形态：普通嵌套下，一个需要 L1 处理的事件要走 L2→L0→L1→L0→L2；HyperTurtle 下，已注册的处理逻辑在 L0 内以 eBPF 形式就地执行，成功后直接恢复 L2；资源不足、竞争或不支持时沿虚线回退到 L1 的原始路径。图中的成功恢复与回退两条分支，与上文 EPT 案例的步数拆解一致——注意 EPT 缺页还包含重执行后再次退出等步骤，成功下放只减少特定处理所需的切换，并不取消 L1 作为正确性与复杂功能兜底的存在。

\subsection{实验证据评议}\label{sec:nested-eval}

两篇论文的硬件、软件年代与实验目标不同，不能把其结果拼成统一的性能排名。Turtles 的宏观负载使用 KVM 同时充当 L0 与 L1，并启用多维分页。相对单层 guest，kernbench 的运行时间增加 14.5\%，SPECjbb 分数下降约 7.8\%；使用修改 L1 的 DRW 配置后，两项数值分别变为 10.3\% 与 6.3\%\cite[\S4.1, Table 2]{source7}。这些差异说明，管理路径密集程度会影响嵌套成本，摘要中的较低开销不能解释为适用于全部工作负载的上界。DRW 还改变了 guest hypervisor，必须与未修改 guest 的基线分别讨论。

HyperTurtle 进一步把热路径收益与端到端效果联系起来。表~\ref{tab:nested-evidence} 只列出各项实验直接报告的比较结果，不将动机分析中的绝对延迟与另一个实验中的加速比拼接成推算值。EPT fault 的平均与尾延迟、服务启动时间、网络往返延迟和采样处理时间是不同度量，不能统一称为“系统加速”。

\begin{table}[tbp]
\centering
\small
\caption[HyperTurtle 的代表性评测结果]{HyperTurtle 的代表性评测结果（原论文报告，非本文复现）。EPT 采用正文\S5.1的数值；图9a的平均改善标注为5.26，二者存在末位差异，不影响此处机制比较。\cite{source8}}
\label{tab:nested-evidence}
\begin{tabularx}{\linewidth}{@{}P{2.4cm}YYY@{}}
\toprule
度量与条件 & 基线 & 报告结果 & 出处 \\
\midrule
EPT fault 处理 & 普通嵌套 L2 & 平均延迟改善约 $5.25\times$，99p 改善 $4.76\times$ & \S5.1，图9a \\
Kata 服务启动 & 相同镜像的普通嵌套 & 4\,KiB 页平均缩短27\%；2\,MiB页缩短8\% & \S5.2，图10 \\
UDP 往返延迟 & Nested-VirtIO & HT+Firewall 平均约60.53\,$\mu$s，基线约90.2\,$\mu$s & \S5.1，表7及正文 \\
采样处理时间，1000\,Hz & L1 对 L2 的常规采样 & HT 对 L2 的处理时间改善 $7.15\times$ & \S5.1，图9b \\
\bottomrule
\end{tabularx}
\end{table}

EPT fault 结果对应的是在 Kata 虚拟机内触发内存按需映射的微基准。减少进入 L1 的往返，使平均延迟接近论文中的非嵌套参考，但只有在负载频繁触发此类缺页时，收益才会显著传导。服务启动实验因此具有解释力：论文从发出启动命令开始计时，直到容器中的 Web 服务能够响应请求，而不是只计 guest 内核 boot。图10还区分了4\,KiB与2\,MiB页；使用大页时，Kata 在 L1 中仍以4\,KiB页映射 L2 内核的部分路径，因而 hyperupcall 仍有作用，但端到端改善较小\cite[\S5.1--5.2]{source8}。这组对照也提醒读者：页大小、后备内存组织和启动阶段都会改变缺页优化的收益。

网络结果须同时比较性能与控制能力。无 L1 策略的 Direct-Assignment 是性能参考，不能冒充功能完全相同的基线。论文用 HT+Pass 与 HT+Firewall 区分基础设施成本和策略程序成本：前者平均延迟约55\,$\mu$s，接近直接分配的54.5\,$\mu$s；后者增加到约60.53\,$\mu$s，但保留了 L1 的防火墙控制，并低于两层 virtio 的约90.2\,$\mu$s\cite[\S5.1]{source8}。由此可以支持“缩短路径而保留策略”的主张，却不能推导出所有 eBPF 网络策略都具有相同开销。profiling 的端到端实验则在相同直接分配网络上改变采样频率：4000\,Hz与1000\,Hz下，论文报告满足500\,$\mu$s查询延迟约束时的 Memcached 吞吐分别提升26\%与12\%\cite[\S5.2, 图14]{source8}。采样更频繁时，减少采样路径切换的价值更明显；这不是无采样负载的通用吞吐提升。

实验外推还受平台与系统配置约束。HyperTurtle 使用双路 Xeon Silver 4216，固定 vCPU 与 vhost 线程在同一 NUMA 节点内的不同物理核，关闭 SMT，并设置\code{idle=poll}与\code{nopti}；吞吐压力实验还增加 Kata 的 vCPU 数量，以使单个 vhost 线程达到瓶颈\cite[\S5, Table 6]{source8}。这些配置改变空闲等待、隔离开销或资源争用条件，不能据此替生产环境预先确定收益。实现规模也应按组成理解：论文 Table 5 分列 L0 KVM、eBPF 框架、QEMU 安装器与动态接口、L1 内存管理和 KVM，以及各个 hyperupcall 的代码量。只引用“KVM 改动151行”会漏掉其他必要部分\cite[\S4.6]{source8}。

\subsection{局限与可检验的后续问题}\label{sec:nested-open}

EPT hyperupcall 的边界在原型中十分具体：它不支持分配大页，对尚未映射进 L1 VMM 页表的部分 page-cache 页面不能直接处理，且全局 EPT hook 暂时只允许一个 L1 为一个 L2 使用该处理器。复杂 L1 功能的并发同步也限制了完整下放，因此原有 L1 路径必须保留\cite[\S4.8]{source8}。这些是特定 EPT 原型的限制，不能扩大成“所有网络 hook 都只能服务一个 L2”；网络路径利用每个虚拟接口的 hook，有不同的实例组织方式。

下一步首先应测量一致性成本，而不只继续压缩最理想路径的延迟。可以逐步增加 L1 与 L2 的并发度，分别控制帧池补充频率、fault log 占用率和锁竞争，报告回退率、99p延迟及宿主 CPU 时间。这样才能判断省下的切换是否被共享状态维护重新消耗。把 EPT hook 扩展为 per-L2 handler 后，还应检验多个 L1 同时注册、撤销与更新程序时的隔离和资源消耗。对生产适用性的评估，则需比较不同 SMT、NUMA、空闲等待和页表隔离设置，而不是假定论文配置代表所有云主机。

安全讨论应区分程序执行约束与系统级信任关系。验证器不能代替 helper 的权限设计、VM归属检查和共享状态协议；可加载程序的执行时间、内存消耗与撤销机制也需要纳入审计。HyperTurtle 假定 L0 能够读取并处理 L1 的相关状态，这与把 L0 排除出信任边界的机密虚拟机模型不同，不能直接移植其安全结论。VM迁移、动态策略更新和更深层嵌套还要求迁移或重新建立程序、maps及hook的关系；这些工作既涉及性能，也涉及生命周期的一致性。

云租户自主管理虚拟机、Kata/serverless 与租户级网络策略为这种方法提供了应用场景。对国内云平台而言，值得评估的是既保留租户控制又减少嵌套管理往返的接口，而不是仅复制一个实验加速比。Turtles 展示了如何把多层语义落实到现有硬件，HyperTurtle 则在这一基础上重新安排部分处理逻辑的位置；其推广价值取决于下放后的状态同步、安全和运维成本能否保持可控。

